% Standalone content fragment. Requires amsmath, amssymb, hyperref.
% Citation keys: dw, richards, bellec, bach, wu, lin.
\section{数学对象与两个不同的风险目标}\label{sec:theory_setup}
设 $X\in\mathbb R^{n\times p}$ 的行为训练特征，$y=X\beta+\varepsilon$。
训练噪声满足 $\mathbb E\varepsilon=0$、$\operatorname{Cov}(\varepsilon)=\sigma^2I_n$，$0<\sigma^2<\infty$，
与系数及设计独立；新特征 $x_0$ 均值为零、协方差为 $\Sigma\succeq0$，
新噪声方差同为 $\sigma^2$，且测试点独立于训练数据。
采用规范化岭损失 $\lVert y-Xb\rVert_2^2/n+\lambda\lVert b\rVert_2^2$，记
\begin{equation}\label{eq:th_ridge}
 S=X^{\mathsf T}X/n,\quad Q_j=(S+\lambda_jI_p)^{-1},\quad
 \widehat\beta_j=Q_jX^{\mathsf T}y/n,\qquad \lambda_j>0.
\end{equation}
候选数为 $K$，混合权重满足 $\boldsymbol1^{\mathsf T}\pi=1$，
估计器为 $\widehat\beta_\pi=\sum_{j=1}^K\pi_j\widehat\beta_j$。

\textbf{给定真实系数和训练结果的风险。}
以 $\mathcal D=(X,y)$ 表示训练数据，定义
\begin{align}
 R_{\mathcal D}(\pi)
 &=\mathbb E[(y_0-x_0^{\mathsf T}\widehat\beta_\pi)^2\mid X,y,\beta]
 =\pi^{\mathsf T}\mathcal K_{\mathcal D}\pi,\label{eq:th_realized_risk}\\
 (\mathcal K_{\mathcal D})_{jk}
 &=\sigma^2+(\widehat\beta_j-\beta)^{\mathsf T}\Sigma(\widehat\beta_k-\beta).
 \label{eq:th_realized_matrix}
\end{align}
它可以用于模拟中的精确总体评估，不需要另抽有限测试集；真实系数仅供评估，不能进入实际选权。
如果以式\eqref{eq:th_realized_matrix}选权，则得到使用真实系数的事后 oracle。

\textbf{给定设计、对系数与训练噪声平均的风险。}
另设 $\mathbb E\beta=0$、$\operatorname{Cov}(\beta)=\Psi\succeq0$，且 $\beta$ 与 $X$ 独立。
对固定权重或仅依赖 $X$ 的权重，定义
$\overline R_X(\pi)=\mathbb E_{\beta,\varepsilon}[R_{\mathcal D}(\pi)\mid X]$。
用此目标选权的 oracle 可知道生成参数 $\Sigma,\Psi,\sigma^2$，但不使用某一次真实系数或噪声的实现值。
两个 oracle 的信息与目标不同，不能混为同一个基准。

\paragraph{命题 1：一般先验下的精确交叉风险。}
在上述二阶矩条件下，
\begin{align}
 \overline R_X(\pi)&=\sigma^2+\pi^{\mathsf T}C_X\pi,\label{eq:th_prior_risk}\\
 (C_X)_{jk}
 &=\lambda_j\lambda_k\operatorname{tr}(\Sigma Q_j\Psi Q_k)
   +\frac{\sigma^2}{n}\operatorname{tr}(\Sigma SQ_jQ_k).
 \label{eq:th_prior_matrix}
\end{align}
矩阵 $C_X$ 对称半正定；式\eqref{eq:th_prior_matrix}不要求 $\Sigma,\Psi,S$ 可交换。

\textit{证明。}
由 $Q_jS=I_p-\lambda_jQ_j$ 和权重和为一，
\begin{equation}\label{eq:th_error_decomp}
 \widehat\beta_\pi-\beta
 =-\left(\sum_j\pi_j\lambda_jQ_j\right)\beta
  +\left(\sum_j\pi_jQ_j\right)X^{\mathsf T}\varepsilon/n.
\end{equation}
两项交叉期望为零。偏差项用 $\mathbb E\beta\beta^{\mathsf T}=\Psi$ 展开，
噪声项用 $\mathbb E\varepsilon\varepsilon^{\mathsf T}=\sigma^2I_n$ 展开，
即得式\eqref{eq:th_prior_matrix}。任意线性组合的对应二次型是两个非负方差项之和，故半正定。
若 $\widehat\pi$ 依赖 $y$，不能把它直接代入固定权重的式\eqref{eq:th_prior_risk}，
并将结果认作 $\mathbb E[R_{\mathcal D}(\widehat\pi)\mid X]$。
\hfill$\square$

\subsection{各向同性情形的差商与高维极限}
当 $\Psi=\alpha^2I_p/p$ 时，令 $t_p(\lambda)=p^{-1}\operatorname{tr}[\Sigma(S+\lambda I_p)^{-1}]$。
先设候选参数两两不同。对 $j\ne k$ 定义
\begin{align}
 B_{jk}&=\frac{t_p(\lambda_j)-t_p(\lambda_k)}{\lambda_k-\lambda_j},&
 V_{jk}&=\frac{\lambda_kt_p(\lambda_k)-\lambda_jt_p(\lambda_j)}{\lambda_k-\lambda_j},\label{eq:th_divided}\\
 B_{jj}&=-t_p'(\lambda_j),&
 V_{jj}&=t_p(\lambda_j)+\lambda_jt_p'(\lambda_j).
 \label{eq:th_diag}
\end{align}
那么
\begin{equation}\label{eq:th_isotropic_matrix}
 (C_X)_{jk}=\alpha^2\lambda_j\lambda_kB_{jk}+\sigma^2(p/n)V_{jk}.
\end{equation}
其证明分别使用
$Q_jQ_k=(Q_j-Q_k)/(\lambda_k-\lambda_j)$ 和
$SQ_jQ_k=(\lambda_kQ_k-\lambda_jQ_j)/(\lambda_k-\lambda_j)$，对角取连续极限。
后者说明方差交叉项包含样本协方差 $S$，不能与偏差交叉项混用。

在 Dobriban--Wager 的比例高维条件下\cite{dw}，例如
$p/n\to\gamma\in(0,\infty)$、$X=Z\Sigma^{1/2}$、$Z$ 元素独立同分布、均值零、方差一且十二阶矩一致有界、
$\lVert\Sigma\rVert_{\mathrm{op}}$ 一致有界、总体经验谱弱收敛到 $H$，
对固定有限候选集、固定 $\lambda_j>0$，可用极限 $t$ 及其导数替换式\eqref{eq:th_divided}--\eqref{eq:th_diag}中的 $t_p$。
其中 $v_\lambda>0$ 解
\begin{equation}\label{eq:th_stieltjes}
 \frac1{v_\lambda}=\lambda+\gamma\int\frac{s}{1+sv_\lambda}\,dH(s),\qquad
 t(\lambda)=\frac1\gamma\left(\frac1{\lambda v_\lambda}-1\right).
\end{equation}
这是现有随机矩阵结论与 resolvent 恒等式的组合，不作为新定理宣称。
一般 $\Psi$ 的偏差项为 $\operatorname{tr}(\Sigma Q_j\Psi Q_k)$，
即使 $\Psi$ 与总体 $\Sigma$ 同时对角化，也通常不与样本 $Q_j$ 交换，因而不能直接使用上述单函数差商。
本项目对一般 $\Psi$ 使用式\eqref{eq:th_prior_matrix}的精确有限样本计算；
一般 source condition 的渐近分析需要另行引用与验证\cite{richards,wu}。

\section{连续最优单岭与平均收益的正确基准}\label{sec:theory_bayes}
\paragraph{命题 2：各向同性二阶矩下的最优线性基准。}
给定 $X$，设 $\Psi=\alpha^2I_p/p$，$\alpha^2,\sigma^2>0$，
令 $V_y=(\alpha^2/p)XX^{\mathsf T}+\sigma^2I_n$。
在全部仅依赖 $X$ 的线性估计器 $Ay$ 中，
\begin{equation}\label{eq:th_lmmse}
 A_*=(\alpha^2/p)X^{\mathsf T}V_y^{-1}
 =(X^{\mathsf T}X+n\lambda_*I_p)^{-1}X^{\mathsf T},\qquad
 \lambda_*=(p/n)\sigma^2/\alpha^2
\end{equation}
具有最小先验平均预测风险，而且
\begin{equation}\label{eq:th_lmmse_gap}
 \overline R_X(A)-\overline R_X(A_*)
 =\operatorname{tr}\!\left[\Sigma(A-A_*)V_y(A-A_*)^{\mathsf T}\right]\ge0.
\end{equation}

\textit{证明。}
$\operatorname{Cov}(\beta-A_*y,y\mid X)=(\alpha^2/p)X^{\mathsf T}-A_*V_y=0$。
将 $Ay-\beta=(A_*y-\beta)+(A-A_*)y$ 的加权平方展开，交叉期望为零，余项即式\eqref{eq:th_lmmse_gap}。
若系数和噪声进一步为独立高斯，$A_*y$ 就是后验均值；
平方损失的条件期望分解将该最优性扩展到全部使用同样 $X,y$ 的估计器，包括依赖 $y$ 的选权。
仅有二阶矩而无高斯假设时，不宣称这个更强结论。
\hfill$\square$

据此，须分别报告平均相对\emph{离散网格最佳单岭}与\emph{连续调参最佳单岭}的差距。
网格未包含 $\lambda_*$ 时，正确先验下平均也可能改善离散网格的最佳值；
先验各向异性后不再有上述普遍的标量岭最优保证，但这并不保证任何给定平均家族都有严格收益。
例如 $S=sI_p$ 时，任意正岭的凸平均仍等价于某个单岭的标量收缩。
使用真实系数的事后 oracle 不属于式\eqref{eq:th_lmmse}的信息范围。

\section{普通聚合 GCV 与交叉残差准则的精确关系}\label{sec:theory_gcv}
记
\begin{equation}\label{eq:th_gcv_notation}
 H_j=XQ_jX^{\mathsf T}/n,\quad r_j=(I_n-H_j)y,\quad
 d_j=1-\operatorname{tr}(H_j)/n,\quad
 \mathsf E=(r_1,\ldots,r_K),\quad D=\operatorname{diag}(d_1,\ldots,d_K).
\end{equation}
正的岭参数保证 $d_j>0$。
以下以 $d=(d_1,\ldots,d_K)^{\mathsf T}$ 表示这些分量构成的向量。
本文所称普通聚合 GCV 与交叉残差矩阵分别为
\begin{align}
 G_{\mathrm{agg}}(\pi)
 &=\frac{\lVert\mathsf E\pi\rVert_2^2/n}{(d^{\mathsf T}\pi)^2},\label{eq:th_agg_gcv}\\
 \widehat{\mathcal K}
 &=D^{-1}\mathsf E^{\mathsf T}\mathsf E D^{-1}/n,\qquad
 G_{\mathrm{cross}}(\pi)=\pi^{\mathsf T}\widehat{\mathcal K}\pi.
 \label{eq:th_cross_gcv}
\end{align}
式\eqref{eq:th_cross_gcv}是既有交叉风险校正框架在完整数据上的形式\cite{bellec}，不作为新方法宣称。
固定权重时，删点并保持惩罚总量 $n\lambda_j$ 不变，逐候选重拟合的 LOO 残差是
\begin{equation}\label{eq:th_loo}
 r_i^{\mathrm{LOO,MA}}=\sum_j\pi_j\frac{r_{j,i}}{1-H_{j,ii}}.
\end{equation}
这是固定权重候选集合的 LOO 准则；若要计算整个自适应选权流程的外层 LOO，还需在删点后重新选权。

\paragraph{命题 3：同一个二次规划可以产生不同的实际模型。}
在单纯形 $\Delta_K=\{\pi\ge0:\boldsymbol1^{\mathsf T}\pi=1\}$ 上，变换
\begin{equation}\label{eq:th_gcv_map}
 u_j=\frac{d_j\pi_j}{d^{\mathsf T}\pi},\qquad
 \pi_j=\frac{u_j/d_j}{\sum_k u_k/d_k}
\end{equation}
是双射，并且
\begin{equation}\label{eq:th_gcv_equiv}
 G_{\mathrm{agg}}(\pi)=u^{\mathsf T}\widehat{\mathcal K}u,\qquad
 \min_{\pi\in\Delta_K}G_{\mathrm{agg}}(\pi)
 =\min_{u\in\Delta_K}G_{\mathrm{cross}}(u).
\end{equation}

\textit{证明。}
$d_j>0$ 保证变换及其逆均将单纯形映到单纯形。
由 $D^{-1}u=\pi/(d^{\mathsf T}\pi)$ 代入式\eqref{eq:th_cross_gcv}，即得等式。
\hfill$\square$

因此，两种优化的最小分数可以完全相同，但普通 GCV 使用的实际权重为 $\pi$，
交叉残差准则使用的实际权重为 $u$。
对正权重分量，$\pi_j/\pi_k=(u_j/u_k)(d_k/d_j)$；
较小 $d_j$ 对应更高自由度，逆变换会提高这些分量的相对比重。
这给出一个无需渐近分析的机制，但并不保证每份数据上普通 GCV 的实际风险都更高。
该变换一般不保持有负权重时的 $\ell_1$ 约束，也可能遇到 $d^{\mathsf T}\pi=0$，故不能原样推广到负权重优化。
这一命题是本项目采用的代数解释；未进行穷尽查新，不据此宣称首次发现。

\subsection{固定权重下普通 GCV 不一致的显式例子}
取独立标准高斯设计，$p/n\to1$，$\Sigma=I_p$，
$\beta\sim N(0,I_p/p)$、$\varepsilon\sim N(0,I_n)$ 相互独立，
固定 $\pi=(1/2,1/2)$。
令 $F$ 为比例一的 Marchenko--Pastur 分布，
$u_0(s)=\sum_j\pi_j\lambda_j/(s+\lambda_j)$。
预测风险与普通 GCV 分别收敛到
\begin{align}
 R_\infty
 &=1+\int\left[u_0(s)^2+\frac{(1-u_0(s))^2}{s}\right]dF(s),\label{eq:th_mp_risk}\\
 G_{\mathrm{agg},\infty}
 &=\frac{\int(s+1)u_0(s)^2\,dF(s)}{\left[\int u_0(s)\,dF(s)\right]^2}.
 \label{eq:th_mp_gcv}
\end{align}
第一式由式\eqref{eq:th_isotropic_matrix}得到；
第二式来自 $y$ 在训练样本谱基中的条件协方差 $XX^{\mathsf T}/p+I_n$。
利用 $t(\lambda)=(\sqrt{1+4/\lambda}-1)/2$，
当 $\lambda=(0.1,1)$ 时二者分别为 $1.787525057$ 和 $1.680756692$；
当 $\lambda=(0.01,10)$ 时分别为 $2.443469136$ 和 $1.769954995$。
这已经是在固定权重下的差异，与额外的数据选权过拟合是不同问题。

\section{受约束选权的稳定性保证与适用范围}\label{sec:theory_stability}
对 $L\ge1$，定义非空紧权重集合
\begin{equation}\label{eq:th_weight_class}
 \mathcal W_L=\{\pi\in\mathbb R^K:\boldsymbol1^{\mathsf T}\pi=1,\ \lVert\pi\rVert_1\le L\}.
\end{equation}
特别地 $\mathcal W_1=\Delta_K$；$L>1$ 允许负权重。

\paragraph{命题 4：矩阵误差到选权误差的确定性传递。}
对任意对称矩阵 $\mathcal K,\widehat{\mathcal K}$，令
$\epsilon=\max_{j,k}|\widehat{\mathcal K}_{jk}-\mathcal K_{jk}|$。
若数值解 $\widehat\pi\in\mathcal W_L$ 满足
$\widehat\pi^{\mathsf T}\widehat{\mathcal K}\widehat\pi
\le\inf_{\pi\in\mathcal W_L}\pi^{\mathsf T}\widehat{\mathcal K}\pi+\delta_{\mathrm{opt}}$，则
\begin{align}
 \sup_{\pi\in\mathcal W_L}
 |\pi^{\mathsf T}(\widehat{\mathcal K}-\mathcal K)\pi|
 &\le L^2\epsilon,\label{eq:th_uniform_bound}\\
 \widehat\pi^{\mathsf T}\mathcal K\widehat\pi
 -\inf_{\pi\in\mathcal W_L}\pi^{\mathsf T}\mathcal K\pi
 &\le 2L^2\epsilon+\delta_{\mathrm{opt}}.
 \label{eq:th_oracle_bound}
\end{align}

\textit{证明。}
逐项取绝对值得
$|\sum_{jk}\pi_j\pi_k(\widehat{\mathcal K}-\mathcal K)_{jk}|
\le\epsilon(\sum_j|\pi_j|)^2\le L^2\epsilon$。
在估计最小值和真实最小值之间加减两个估计二次型，分别用此界与优化误差，即得第二式。
\hfill$\square$

这个界本身不说明 $\epsilon\to0$。它说明无约束放大负权重会同步放大风险估计误差；
$L=3$ 的界常数是 $L=1$ 时的九倍，而非免费的模型类扩张。

\paragraph{推论 1：已有交叉风险界可支持的固定候选集结论。}
考虑一列高斯线性回归问题：$x_i\sim N(0,\Sigma_n)$，独立噪声为 $N(0,\sigma^2)$，
$\sup_n p/n<\infty$、$\sup_n\lVert\Sigma_n\rVert_{\mathrm{op}}<\infty$；
固定 $K,L$，确定性候选参数满足
$0<\underline\lambda\le\lambda_j\le\overline\lambda<\infty$。
令固定系数满足 $\sup_n\beta_n^{\mathsf T}\Sigma_n\beta_n<\infty$。
则式\eqref{eq:th_cross_gcv}的估计矩阵对式\eqref{eq:th_realized_matrix}满足
\begin{equation}\label{eq:th_cross_rate}
 \max_{j,k}|\widehat{\mathcal K}_{jk}-(\mathcal K_{\mathcal D})_{jk}|
 =O_{\mathbb P}(n^{-1/2}),
\end{equation}
因此在数值优化误差为 $o_{\mathbb P}(1)$ 时，
\begin{equation}\label{eq:th_cross_consistency}
 R_{\mathcal D}(\widehat\pi)-\inf_{\pi\in\mathcal W_L}R_{\mathcal D}(\pi)
 =O_{\mathbb P}(L^2n^{-1/2})+\delta_{\mathrm{opt}}=o_{\mathbb P}(1).
\end{equation}
这里的 oracle 目标是给定真实系数与训练结果的 $R_{\mathcal D}$，不是先验平均 $\overline R_X$。

\textit{推导与引用边界。}
使用 Bellec 等人交叉风险非渐近界\cite{bellec}的 Theorem 1、式 (20)，
取完整数据 $I_j=[n]$；这个逐对界允许不同的惩罚 $g_j(b)=\lambda_j\lVert b\rVert_2^2/2$。
在上述条件下，它给出
\begin{equation}\label{eq:th_cited_cross_bound}
 \mathbb E\left[
 d_jd_k\frac{|\widehat{\mathcal K}_{jk}-(\mathcal K_{\mathcal D})_{jk}|}
 {\sqrt{(\mathcal K_{\mathcal D})_{jj}(\mathcal K_{\mathcal D})_{kk}}}
 \right]\le\frac{C_0}{\sqrt n},
\end{equation}
其中 $C_0$ 可依赖固定的比例、谱和正则界。高斯设计给出 $\lVert S\rVert_{\mathrm{op}}=O_{\mathbb P}(1)$，
而 $d_j\ge\underline\lambda/(\lVert S\rVert_{\mathrm{op}}+\underline\lambda)$，故所有 $d_j$ 以高概率远离零。
由岭线性算子的奇异值界，
$\lVert\widehat\beta_j\rVert_2\le\lVert y\rVert_2/(2\sqrt{n\underline\lambda})$；
再用 $\lVert y\rVert_2^2/n=O_{\mathbb P}(1)$ 可得 $(\mathcal K_{\mathcal D})_{jj}=O_{\mathbb P}(1)$。
式\eqref{eq:th_cited_cross_bound}、Markov 不等式与固定有限个指标的并界得到式\eqref{eq:th_cross_rate}；
最后应用命题 4。
对独立随机系数，只要 $\beta^{\mathsf T}\Sigma_n\beta=O_{\mathbb P}(1)$，同样可条件化使用这个论证。

这是一条由现有逐对估计界与确定性误差传递得到的受限推论，不是对任意候选数、任意负权重或 $\lambda\to0$ 的完整一致性定理。
尤其不能把文献中随机子样本、共同惩罚的普通 GCV 不一致定理，原样当作同一设计不同 $\lambda$ 的定理。
式\eqref{eq:th_cited_cross_bound}的逐对结果足以用于这里的固定候选集，但须核查所用完整数据矩阵和参数界。

\section{信号谱对齐、滤波器限制与本次结论的范围}\label{sec:theory_source}
设总体协方差的特征对为 $(\mu_i,v_i)$，$\mu_i\asymp i^{-a}$，$a>1$。
参照 source condition，令
\begin{equation}\label{eq:th_source_condition}
 \mathbb E\langle v_i,\beta\rangle\langle v_j,\beta\rangle=0\ (i\ne j),\qquad
 \mu_i\mathbb E\langle v_i,\beta\rangle^2\asymp i^{-b}.
\end{equation}
这里是对先验的二阶矩条件，不是某一次随机系数的逐坐标等式\cite{lin,richards}。
相应系数方差与 $i^{a-b}$ 成比例；$b=a$ 对应各向同性，
$b<a$ 相对提高弱谱方向的系数方差，$b>a$ 相对强调强谱方向。
比较 $b$ 时固定 $\operatorname{tr}(\Sigma\Psi)/\sigma^2$，才能排除总预测信号能量变化的影响。
响应机制仍为正确线性模型，故应称为先验或正则结构变化，而非自动称线性模型误设。

在样本谱值 $s\ge0$ 上，平均岭的拟合滤波器为
\begin{equation}\label{eq:th_filter}
 h_\pi(s)=\sum_j\pi_j\frac{s}{s+\lambda_j}.
\end{equation}
凸权重下 $h_\pi'(s)>0$、$h_\pi''(s)<0$；允许负权重可解除这些形状限制，
但固定有限 $K$ 时仍只是带固定极点的有理函数族，满足 $h_\pi(0)=0$、$h_\pi(\infty)=1$。
它不能恢复训练设计零空间中的系数，扩大滤波器家族也不等于自动获得可实现的低风险选权。
跨岭参数的外推与加权正则化已有直接研究\cite{bach,wu}，本报告不把这些组合本身称为首次提出。

\textbf{有限维实验与渐近保证的区别。}
若将截断幂律谱重新规范到 $\operatorname{tr}(\Sigma)=p$，其最大特征值会随 $p$ 增长，
通常不满足上述一致有界谱假设；若不重标，经验谱又可能退化。
因此本次有限维机制实验的精确风险矩阵始终有效，但它不能单独证明无限维 Scaling Law，
也不能将推论 1 的一致性保证无条件外推到该幂律增维路径。
本次可争取的结果是：确认一般先验的有限样本风险基线，分开测量网格收益、模型类收益和实际选权代价，
并在满足明确假设的固定候选集上陈述可证明的稳定性结论。
